\chapter{Pfaffian и \texorpdfstring{\(\eta\)}{eta}-phase: reduced sign и KO6 cancellation}

Bosonic spectral action не различает две orientation branches. Проверим,
содержит ли fermionic measure дополнительную фазовую информацию.

\section{Каноническая chiral antisymmetric form}

В plus--minus ordering bipartite primary operator имеет вид
\begin{equation}
 D_{\rm red}=
 \begin{pmatrix}
 0&M\\
 M^T&0
 \end{pmatrix},
 \qquad
 \gamma=
 \begin{pmatrix}
 I_6&0\\0&-I_6
 \end{pmatrix}.
\end{equation}
Поэтому
\begin{equation}
 \mathcal A_{\rm red}:=\gamma D_{\rm red}
 =
 \begin{pmatrix}
 0&M\\
 -M^T&0
 \end{pmatrix}
\end{equation}
антисимметрична и допускает Pfaffian. При выбранной orientation Grassmann
basis
\begin{equation}
 \operatorname{Pf}\mathcal A_{\rm red}
 =
 (-1)^{6\cdot5/2}\det M
 =-\det M.
\end{equation}

Для двух minima точный расчёт даёт
\begin{center}
\begin{tabular}{lcc}
\toprule
Ветвь & \(\det M\) & \(\operatorname{Pf}\mathcal A_{\rm red}\)\\
\midrule
\(\mathcal B_-\) & \(+3/2\) & \(-3/2\)\\
\(\mathcal B_+\) & \(-3/2\) & \(+3/2\)\\
\bottomrule
\end{tabular}
\end{center}

\begin{proposition}[Conditional reduced Pfaffian phase]
На одном ориентированном reduced representative две branches различаются
relative Pfaffian phase \(\pi\), хотя их bosonic spectra совпадают.
\end{proposition}

\section{Полное KO6-сокращение}

KO6 completion добавляет \(J\)-сопряжённую копию с противоположной grading.
Полная antisymmetric form имеет block structure
\begin{equation}
 \mathcal A_{\rm full}
 =
 \mathcal A_{\rm red}\oplus(-\overline{\mathcal A}_{\rm red}).
\end{equation}
Поскольку reduced block имеет dimension \(12=2\cdot6\),
\begin{equation}
 \operatorname{Pf}\mathcal A_{\rm full}
 =
 \left|\operatorname{Pf}\mathcal A_{\rm red}\right|^2
 =\frac94
\end{equation}
в обеих branches. Полный \(J\)-paired functional measure теряет relative
phase.

\begin{theorem}[Full KO6 Pfaffian cancellation]
Если интегрировать обе \(J\)-сопряжённые copies как независимые Grassmann
variables, Pfaffian orientation sign сокращается. Selector может сохраниться
только в physical reality-orbit measure, считающей одну copy, а не в полном
doubled Pfaffian.
\end{theorem}

\section{Почему reduced sign ещё не является предсказанием}

Нечётная перестановка Grassmann basis меняет знак reduced Pfaffian. Поэтому
для физического relative sign нужна глобально заданная orientation
determinant line, которой текущий finite graph не выводит.

Кроме того, при непрерывном complex connector \(z=e^{i\theta}\)
\begin{align}
 \det M_-&=\frac12(z+2)z^{-3},\\
 \det M_+&=-\frac12(z-2)z^{-3}.
\end{align}
На единичной окружности эти determinants не обращаются в нуль. Их phases
могут непрерывно вращаться, поэтому real signs в minima не являются
\(\mathbb Z_2\)-индексом полного complex configuration space.

Наконец, chiral symmetry даёт
\begin{equation}
 \eta(D_-)=\eta(D_+)=0.
\end{equation}
Pfaffian sign создаёт quantum relative sign, но не разность вещественных
vacuum energies и сам по себе не выбирает одну classical branch.

\begin{nogocriterion}[Запрет неориентированного Pfaffian selector]
Нельзя считать branch выбранной только потому, что reduced Pfaffians имеют
разные знаки. Требуется вывести orientation determinant line и показать,
почему physical measure использует один reality representative, тогда как
\(J\)-copy обеспечивает consistency, но не умножает Pfaffian обратно.
\end{nogocriterion}

Итак, Pfaffian route не закрыт отрицательно: он впервые обнаруживает relative
phase \(\pi\). Но phase условна и сокращается в полном KO6 measure.
Следующий гейт --- determinant-line/anomaly-inflow construction, способная
сделать reduced sign basis-independent. Exact ledger сохранён в
\path{s2t_v4_pfaffian_eta_orientation_gate_results.json}.